\documentclass[11pt,a4paper]{article}

\usepackage[utf8]{inputenc}
\usepackage[T1]{fontenc}
\usepackage[provide=*,english]{babel}
\usepackage{lmodern}
\usepackage{microtype}
\usepackage{amsmath,amssymb,mathtools}
\usepackage{geometry}
\usepackage{xcolor}
\usepackage{enumitem}
\usepackage{booktabs}
\usepackage{array}
\usepackage{fancyhdr}
\usepackage[hidelinks]{hyperref}

\geometry{left=2.3cm,right=2.3cm,top=2.2cm,bottom=2.25cm,headheight=14pt}
\setlength{\parindent}{0pt}
\setlength{\parskip}{0.38em}
\setlist[itemize]{topsep=0.25em,itemsep=0.18em,leftmargin=1.5em}
\setlist[enumerate]{topsep=0.25em,itemsep=0.22em,leftmargin=1.8em}
\allowdisplaybreaks

\definecolor{blue}{RGB}{34,73,110}
\definecolor{gray}{RGB}{90,96,102}

\pagestyle{fancy}
\fancyhf{}
\fancyhead[L]{\small\color{gray}GC2D: ABBA method}
\fancyhead[R]{\small\color{gray}ABBA2 implicit: reduced multiplier solve}
\fancyfoot[C]{\thepage}
\renewcommand{\headrulewidth}{0.3pt}

\newcommand{\R}{\mathbb{R}}
\newcommand{\Id}{I}
\newcommand{\Y}{\mathcal{Y}}
\newcommand{\M}{\mathcal{M}}
\newcommand{\Om}{\Omega_{\times}}
\newcommand{\PhiAB}{\Phi^{\mathrm{ABBA}}_{h,t_n}}
\newcommand{\norm}[1]{\left\lVert #1\right\rVert}

\begin{document}

\begin{center}
  {\LARGE\bfseries\color{blue}ABBA2 Implicit with Hairer's Symmetric Projection}\par
  \vspace{0.35em}
  {\large Reduced physical derivation and canonical configuration space}\par
  \vspace{0.45em}
  {\small Endpoint-time $A$--$B$--$B$--$A$ map with symmetric diagonal projection}
\end{center}

\vspace{0.5em}

\section{Scope and implementation name}

The public ABBA family contains exactly five numerical-method classes:
\begin{itemize}
 \item \texttt{ABBA2Midpoint};
 \item \texttt{ABBA2Implicit};
 \item \texttt{ABBA4Implicit};
 \item \texttt{ABBA4ImplicitSingleProjection};
 \item \texttt{ABBA6Implicit}.
\end{itemize}
State-space variants are parameters of these classes, not additional public
method classes.  The four implicit methods expose three independent axes:

\begin{center}
\small
\begin{tabular}{@{}>{\raggedright\arraybackslash}p{0.27\textwidth}
                    >{\raggedright\arraybackslash}p{0.65\textwidth}@{}}
\toprule
Axis & Canonical values \\
\midrule
\texttt{projection\_formulation}
  & \texttt{reduced\_multiplier};
    \texttt{simultaneous\_state\_}\allowbreak\texttt{multiplier} \\
\texttt{nonlinear\_solver}
  & \texttt{newton}; \texttt{broyden} \\
\texttt{state\_extension}
  & \texttt{physical}; \texttt{shared\_time}; \texttt{fully\_extended} \\
\bottomrule
\end{tabular}
\end{center}

Each implicit class therefore admits $2\times2\times3=12$ configurations.
\texttt{ABBA2Midpoint} has no nonlinear residual and accepts only the three
state extensions.  The five public classes consequently represent
$4\times12+3=51$ valid configurations.

The formulation and solver choices are global within a composed step.  All
three signed maps of \texttt{ABBA4Implicit}, and all seven signed maps of
\texttt{ABBA6Implicit}, use the same configuration.  In contrast,
\texttt{ABBA4ImplicitSingleProjection} surrounds its complete unprojected
triple jump with one projection and therefore performs one nonlinear solve per
outer step.

This note derives the following branch of \texttt{ABBA2Implicit}:
\begin{itemize}
 \item \texttt{projection\_formulation}: \texttt{reduced\_multiplier};
 \item \texttt{state\_extension}: \texttt{physical}.
\end{itemize}
Its nonlinear unknown is $\mu\in\R^2$ per particle and its exact Newton matrix
is $2\times2$.  The simultaneous formulation instead solves for
$(u_f,v_f,\mu)\in\R^6$ per particle with a $6\times6$ Newton block.  Both
formulations define the same exact projected map when converged; Newton and
Broyden are alternative algorithms for finding that root.  The
six-dimensional simultaneous vector is only an algebraic solver workspace,
not a trajectory state or a separate numerical method.

\section{Guiding-center system and duplicated phase space}

The physical state is
\[
 z=(X,Y)\in\R^2,
\]
and, with the convention used here, it satisfies
\[
 \dot X=-\partial_Y\phi(t,X,Y),
 \qquad
 \dot Y=\partial_X\phi(t,X,Y).
\]
We define
\[
 J_0=
 \begin{pmatrix}
  0&-1\\
  1&0
 \end{pmatrix},
 \qquad
 f(z,t)=J_0\nabla_z\phi(z,t),
\]
so that $\dot z=f(z,t)$. Its spatial Jacobian is
\[
 W(z,t):=D_zf(z,t)=J_0\nabla_z^2\phi(z,t).
\]
Since the Hessian of $\phi$ is symmetric,
\[
 W(z,t)^TJ_0+J_0W(z,t)=0. \tag{1}
\]

We duplicate the state:
\[
 \Y=(u,v)=(X_1,Y_1,X_2,Y_2)\in\R^4,
 \qquad u,v\in\R^2.
\]
The duplicated dynamics associated with the subflows used below is
\[
 \dot u=f(v,t),
 \qquad
 \dot v=f(u,t).
\]
If initially $u=v=z$, both equations coincide with the physical dynamics.

The physical state is identified with the diagonal manifold
\[
 \M=\{(u,v)\in\R^4:u=v\}.
\]
We write it as
\[
 g(\Y)=G\Y=u-v,
 \qquad
 G=(\Id_2\; -\Id_2),
 \qquad
 GG^T=2\Id_2.
\]
For $\mu=(\mu_X,\mu_Y)^T$,
\[
 G^T\mu=
 \begin{pmatrix}\mu\\-\mu\end{pmatrix}
 =(\mu_X,\mu_Y,-\mu_X,-\mu_Y)^T.
\]

The constraints are linear. Therefore,
\[
 \nabla^2g_1=\nabla^2g_2=0
\]
and, for every $\mu$,
\[
 \boxed{
 H_\mu(\Y)=\sum_{j=1}^2\mu_j\nabla^2g_j(\Y)=0.
 }
\]
This does not mean that $\mu=0$: $H_\mu$ is identically zero because $G=Dg$ is constant.

\section{The explicit submaps \texorpdfstring{$A$}{A} and \texorpdfstring{$B$}{B}}

For a substep length $s$ and an evaluation time $t$, we define
\[
 A_{s,t}(u,v)=\bigl(u+s f(v,t),v\bigr), \tag{2}
\]
\[
 B_{s,t}(u,v)=\bigl(u,v+s f(u,t)\bigr). \tag{3}
\]
The map $A$ advances only the first copy using the second one; $B$ advances only the
second copy using the first one. Both maps are explicit.

Their Jacobians are
\[
 DA_{s,t}(u,v)=
 \begin{pmatrix}
  \Id_2&sW(v,t)\\
  0&\Id_2
 \end{pmatrix},
 \qquad
 DB_{s,t}(u,v)=
 \begin{pmatrix}
  \Id_2&0\\
  sW(u,t)&\Id_2
 \end{pmatrix}. \tag{4}
\]

In the duplicated phase space, we use the skew-symmetric matrix
\[
 \Om=
 \begin{pmatrix}
  0&J_0\\
  J_0&0
 \end{pmatrix}.
\]
Identity (1) directly implies
\[
 (DA_{s,t})^T\Om DA_{s,t}=\Om,
 \qquad
 (DB_{s,t})^T\Om DB_{s,t}=\Om.
\]
Thus, each submap is symplectic with respect to $\Om$, and any composition of these maps is
also symplectic in the duplicated phase space.

Moreover, the variables used to evaluate the field remain fixed in each submap. Therefore,
\[
 \boxed{A_{s,t}^{-1}=A_{-s,t},\qquad B_{s,t}^{-1}=B_{-s,t}.} \tag{5}
\]

\section{The nonautonomous \texorpdfstring{$A$--$B$--$B$--$A$}{A-B-B-A} step}

Let
\[
 t_{n+1}=t_n+h,
 \qquad
 s=\frac h2.
\]
The base map that replaces BM4 is
\[
 \boxed{
 \PhiAB
 =A_{s,t_{n+1}}
  \circ B_{s,t_{n+1}}
  \circ B_{s,t_n}
  \circ A_{s,t_n}.
 } \tag{6}
\]
Compositions are read from right to left. Therefore, the effective order is
\[
 A_{s,t_n}\longrightarrow B_{s,t_n}
 \longrightarrow B_{s,t_{n+1}}\longrightarrow A_{s,t_{n+1}}.
\]

Starting from $(u^{(0)},v^{(0)})$, the four stages are
\begin{align}
 u^{(1)}
 &=u^{(0)}+s f\bigl(v^{(0)},t_n\bigr), &&(A_1) \tag{7a}\\
 v^{(1)}
 &=v^{(0)}+s f\bigl(u^{(1)},t_n\bigr), &&(B_1) \tag{7b}\\
 v^{(2)}
 &=v^{(1)}+s f\bigl(u^{(1)},t_{n+1}\bigr), &&(B_2) \tag{7c}\\
 u^{(2)}
 &=u^{(1)}+s f\bigl(v^{(2)},t_{n+1}\bigr). &&(A_2) \tag{7d}
\end{align}
The output is
\[
 \PhiAB(u^{(0)},v^{(0)})=(u^{(2)},v^{(2)}).
\]

The evaluations at the two time endpoints are essential in the nonautonomous problem.
If every stage were evaluated at $t_n$, time symmetry would be lost. In the autonomous
case, the two central stages can be combined, and one recovers
\[
 A_{h/2}\longrightarrow B_h\longrightarrow A_{h/2}.
\]

Since $A$ and $B$ are symplectic with respect to $\Om$, the complete map is also symplectic:
\[
 \boxed{
 (D\PhiAB)^T\Om D\PhiAB=\Om.
 } \tag{8}
\]

\section{Proof of the symmetry of ABBA}

A time-dependent method is symmetric if
\[
 \Phi^{\mathrm{ABBA}}_{-h,t_{n+1}}
 =\left(\Phi^{\mathrm{ABBA}}_{h,t_n}\right)^{-1}. \tag{9}
\]
Using (5) and reversing the order of composition in (6),
\begin{align*}
 \left(\Phi^{\mathrm{ABBA}}_{h,t_n}\right)^{-1}
 &={}
 A_{-s,t_n}
 \circ B_{-s,t_n}
 \circ B_{-s,t_{n+1}}
 \circ A_{-s,t_{n+1}}.
\end{align*}
We now start at $t_{n+1}$ with a step of size $-h$. Its final time is
$t_{n+1}-h=t_n$, and definition (6) gives precisely
\[
 \Phi^{\mathrm{ABBA}}_{-h,t_{n+1}}
 =A_{-s,t_n}
  \circ B_{-s,t_n}
  \circ B_{-s,t_{n+1}}
  \circ A_{-s,t_{n+1}}.
\]
Therefore,
\[
 \boxed{
 \Phi^{\mathrm{ABBA}}_{-h,t_{n+1}}
 =\left(\Phi^{\mathrm{ABBA}}_{h,t_n}\right)^{-1}.
 }
\]
The architecture is palindromic when both the sign of the step and the two endpoint times
are reversed simultaneously.

\section{Hairer's symmetric projection around ABBA}

Suppose that
\[
 \Y_n=(z_n,z_n)\in\M
\]
is known. For a multiplier $\mu\in\R^2$, we first correct the input:
\[
 \widehat{\Y}_n(\mu)
 =\Y_n+G^T\mu
 =(z_n+\mu,z_n-\mu).
\]
We write
\[
 u^{(0)}=z_n+\mu,
 \qquad
 v^{(0)}=z_n-\mu,
\]
and apply the four stages in (7). The auxiliary result is
\[
 \widetilde{\Y}_{n+1}(\mu)
 =\PhiAB\bigl(\Y_n+G^T\mu\bigr)
 =(u^{(2)}(\mu),v^{(2)}(\mu)).
\]
The same correction is then applied at the output:
\[
 \Y_{n+1}(\mu)
 =\widetilde{\Y}_{n+1}(\mu)+G^T\mu
 =\bigl(u^{(2)}+\mu,v^{(2)}-\mu\bigr). \tag{10}
\]
Imposing $\Y_{n+1}\in\M$ is equivalent to
\[
 u^{(2)}(\mu)+\mu=v^{(2)}(\mu)-\mu.
\]
Thus, the problem is reduced exactly to two equations:
\[
 \boxed{
 R(\mu)
 :=u^{(2)}(\mu)-v^{(2)}(\mu)+2\mu=0.
 } \tag{11}
\]
In matrix notation,
\[
 R(\mu)
 =G\PhiAB\bigl(\Y_n+G^T\mu\bigr)+2\mu. \tag{12}
\]
Once the root $\mu_*$ has been found, the new physical state is
\[
 \boxed{
 z_{n+1}=u^{(2)}(\mu_*)+\mu_*
          =v^{(2)}(\mu_*)-\mu_*.
 } \tag{13}
\]

\section{Exact Jacobian of the ABBA map}

For a particular evaluation of (7), we define
\begin{align*}
 W_1&=W\bigl(v^{(0)},t_n\bigr),
 &W_2&=W\bigl(u^{(1)},t_n\bigr),\\
 W_3&=W\bigl(u^{(1)},t_{n+1}\bigr),
 &W_4&=W\bigl(v^{(2)},t_{n+1}\bigr),
\end{align*}
and abbreviate
\[
 S:=W_2+W_3.
\]

\subsection*{Explicit expansion of \texorpdfstring{$u^{(2)}$}{u(2)}}

To compute the derivative of the last stage without hiding any dependence, we start from
\[
 u^{(2)}=u^{(1)}+s f\bigl(v^{(2)},t_{n+1}\bigr)
\]
and successively substitute the preceding stages. First,
\begin{align*}
 v^{(2)}
 &=v^{(1)}+s f\bigl(u^{(1)},t_{n+1}\bigr)\\
 &=v^{(0)}+s f\bigl(u^{(1)},t_n\bigr)
   +s f\bigl(u^{(1)},t_{n+1}\bigr).
\end{align*}
Since
\[
 u^{(1)}=u^{(0)}+s f\bigl(v^{(0)},t_n\bigr),
\]
the expression for $v^{(2)}$ solely in terms of the initial data is
\[
 \begin{aligned}
 v^{(2)}={}&v^{(0)}
 +s f\!\left(u^{(0)}+s f\bigl(v^{(0)},t_n\bigr),t_n\right)\\
 &+s f\!\left(u^{(0)}+s f\bigl(v^{(0)},t_n\bigr),t_{n+1}\right).
 \end{aligned}
\]
Therefore, complete substitution into the last stage gives
\[
 \boxed{
 \begin{aligned}
 u^{(2)}={}&u^{(0)}+s f\bigl(v^{(0)},t_n\bigr)\\
 &+s f\!\Bigl(
 v^{(0)}
 +s f\!\left(u^{(0)}+s f\bigl(v^{(0)},t_n\bigr),t_n\right)\\
 &\hspace{2.7cm}
 +s f\!\left(u^{(0)}+s f\bigl(v^{(0)},t_n\bigr),t_{n+1}\right),
 t_{n+1}\Bigr).
 \end{aligned}
 }
\]
This formula displays directly all dependencies of $u^{(2)}$ on $u^{(0)}$ and $v^{(0)}$.
Nevertheless, to apply the chain rule readably, it is convenient to differentiate the
intermediate stages in the same order in which they are computed.

For the first stage, we have
\[
 \frac{\partial u^{(1)}}{\partial u^{(0)}}=\Id_2,
 \qquad
 \frac{\partial u^{(1)}}{\partial v^{(0)}}=sW_1.
\]
Differentiating the expression for $v^{(2)}$ yields
\begin{align*}
 \frac{\partial v^{(2)}}{\partial u^{(0)}}
 &=sW_2\frac{\partial u^{(1)}}{\partial u^{(0)}}
   +sW_3\frac{\partial u^{(1)}}{\partial u^{(0)}}
 =s(W_2+W_3)=sS,\\
 \frac{\partial v^{(2)}}{\partial v^{(0)}}
 &=\Id_2
   +sW_2\frac{\partial u^{(1)}}{\partial v^{(0)}}
   +sW_3\frac{\partial u^{(1)}}{\partial v^{(0)}}\\
 &=\Id_2+s^2(W_2+W_3)W_1
 =\Id_2+s^2SW_1.
\end{align*}
Finally, differentiating
$u^{(2)}=u^{(1)}+s f(v^{(2)},t_{n+1})$, we obtain
\begin{align*}
 \frac{\partial u^{(2)}}{\partial u^{(0)}}
 &=\frac{\partial u^{(1)}}{\partial u^{(0)}}
   +sW_4\frac{\partial v^{(2)}}{\partial u^{(0)}}\\
 &=\Id_2+s^2W_4S,\\[1mm]
 \frac{\partial u^{(2)}}{\partial v^{(0)}}
 &=\frac{\partial u^{(1)}}{\partial v^{(0)}}
   +sW_4\frac{\partial v^{(2)}}{\partial v^{(0)}}\\
 &=sW_1+sW_4\bigl(\Id_2+s^2SW_1\bigr)\\
 &=s(W_1+W_4)+s^3W_4SW_1.
\end{align*}
Consequently, the two blocks in the first row of the ABBA Jacobian are
\[
 \boxed{
 A:=\frac{\partial u^{(2)}}{\partial u^{(0)}}
 =\Id_2+s^2W_4S,
 \qquad
 B:=\frac{\partial u^{(2)}}{\partial v^{(0)}}
 =s(W_1+W_4)+s^3W_4SW_1.
 }
\]
The order of the matrix products follows from the chain rule and must be preserved because
the matrices $W_i$ evaluated at different points need not commute.

The Jacobians of the four stages, evaluated at the points actually traversed, are
\[
 M_{A_1}=
 \begin{pmatrix}
  \Id_2&sW_1\\0&\Id_2
 \end{pmatrix},
 \qquad
 M_{B_1}=
 \begin{pmatrix}
  \Id_2&0\\sW_2&\Id_2
 \end{pmatrix},
\]
\[
 M_{B_2}=
 \begin{pmatrix}
  \Id_2&0\\sW_3&\Id_2
 \end{pmatrix},
 \qquad
 M_{A_2}=
 \begin{pmatrix}
  \Id_2&sW_4\\0&\Id_2
 \end{pmatrix}.
\]
By the chain rule,
\[
 \boxed{
 D\PhiAB=M_{A_2}M_{B_2}M_{B_1}M_{A_1}.
 } \tag{14}
\]
The two central matrices satisfy
\[
 M_{B_2}M_{B_1}
 =
 \begin{pmatrix}
  \Id_2&0\\sS&\Id_2
 \end{pmatrix}.
\]
Multiplying the four blocks gives the exact $4\times4$ Jacobian:
\[
 \boxed{
 D\PhiAB=
 \begin{pmatrix}
  \Id_2+s^2W_4S
  &s(W_1+W_4)+s^3W_4SW_1\\[2mm]
  sS
  &\Id_2+s^2SW_1
 \end{pmatrix}.
 } \tag{15}
\]
The Jacobian of the map has not been replaced by the identity: (15) contains exactly all
dependencies generated by the four stages.

\section{Exact Jacobian of the reduced residual}

From
\[
 \frac{\partial(u^{(0)},v^{(0)})}{\partial\mu}
 =\begin{pmatrix}\Id_2\\-\Id_2\end{pmatrix}=G^T
\]
and from (12), we obtain
\[
 \boxed{
 R'(\mu)=G\,D\PhiAB\bigl(\Y_n+G^T\mu\bigr)G^T+2\Id_2.
 } \tag{16}
\]
If we write (15) as
\[
 D\PhiAB=
 \begin{pmatrix}A&B\\C&D\end{pmatrix},
\]
then
\[
 G\,D\PhiAB\,G^T=A-B-C+D.
\]
Substituting the four blocks in (15) yields
\[
 \boxed{
 \begin{aligned}
 R'(\mu)={}&4\Id_2
 -s\bigl(W_1+W_2+W_3+W_4\bigr)\\
 &+s^2\bigl[W_4(W_2+W_3)+(W_2+W_3)W_1\bigr]\\
 &-s^3W_4(W_2+W_3)W_1.
 \end{aligned}
 } \tag{17}
\]
Since $s=h/2$, this can also be written as
\[
 \boxed{
 \begin{aligned}
 R'(\mu)={}&4\Id_2
 -\frac h2\bigl(W_1+W_2+W_3+W_4\bigr)\\
 &+\frac{h^2}{4}
 \bigl[W_4(W_2+W_3)+(W_2+W_3)W_1\bigr]\\
 &-\frac{h^3}{8}W_4(W_2+W_3)W_1.
 \end{aligned}
 } \tag{18}
\]
The matrices $W_i$ need not commute: the order of the products in (17) and (18) must be
preserved exactly.

\section{Explicit Jacobian of the implicit physical step}

The Jacobian $R'(\mu)$ computed in the previous section is the matrix used to solve the
nonlinear system that determines the multiplier. However, it is not the Jacobian of the
physical numerical method. The object needed, for example, to verify the symplecticity of
the step directly is
\[
 D\Psi_{h,t_n}(z_n)
 =\frac{\partial z_{n+1}}{\partial z_n}\in\R^{2\times2},
\]
where
\[
 \Psi_{h,t_n}:z_n\longmapsto z_{n+1}.
\]
The difference between these matrices is essential: changing $z_n$ also changes the
converged root $\mu_*=\mu_h(z_n)$. Therefore, computing $D\Psi_{h,t_n}(z_n)$ requires
implicit differentiation of the multiplier.

Throughout this section, the matrices $W_i$ and the Jacobian $D\PhiAB$ are evaluated along
the ABBA stages corresponding to the converged solution $\mu_h(z_n)$. In particular, for
GC2D,
\[
 W_i=
 \begin{pmatrix}
  -\phi_{XY}&-\phi_{YY}\\
  \phi_{XX}&\phi_{XY}
 \end{pmatrix}_{(z_i,t_i)}.
\]
The four points $(z_i,t_i)$ are, respectively,
$(v^{(0)},t_n)$, $(u^{(1)},t_n)$, $(u^{(1)},t_{n+1})$, and
$(v^{(2)},t_{n+1})$.

\subsection{Partial derivatives of the residual}

We introduce the matrices
\[
 E=
 \begin{pmatrix}\Id_2\\\Id_2\end{pmatrix},
 \qquad
 N=G^T=
 \begin{pmatrix}\Id_2\\-\Id_2\end{pmatrix}.
\]
The ABBA input state can be written as
\[
 \begin{pmatrix}u^{(0)}\\v^{(0)}\end{pmatrix}
 =Ez_n+N\mu
 =\begin{pmatrix}z_n+\mu\\z_n-\mu\end{pmatrix}.
\]
To display the dependence on the initial state, we write the residual as a function of its
two variables:
\[
 \mathcal R(z,\mu)
 =G\PhiAB(Ez+N\mu)+2\mu.
\]
Let
\[
 \mathcal J_{\mathrm{ABBA}}
 :=D\PhiAB
 =\begin{pmatrix}A&B\\C&D\end{pmatrix},
\]
with the blocks from (15), namely,
\[
 \begin{aligned}
 A&=\Id_2+s^2W_4S,
 &B&=s(W_1+W_4)+s^3W_4SW_1,\\
 C&=sS,
 &D&=\Id_2+s^2SW_1.
 \end{aligned}
\]

When taking a partial derivative with respect to $\mu$, the state $z$ is held fixed. Since
\[
 \frac{\partial(Ez+N\mu)}{\partial\mu}=N,
\]
the chain rule gives
\[
 \boxed{
 \mathcal K
 :=D_\mu\mathcal R
 =G\mathcal J_{\mathrm{ABBA}}N+2\Id_2.
 }
\]
Similarly, when differentiating with respect to $z$, $\mu$ is held fixed and
\[
 \frac{\partial(Ez+N\mu)}{\partial z}=E.
\]
Therefore,
\[
 \boxed{
 \mathcal L
 :=D_z\mathcal R
 =G\mathcal J_{\mathrm{ABBA}}E.
 }
\]
In terms of the four blocks,
\[
 \boxed{
 \mathcal K=A-B-C+D+2\Id_2,
 \qquad
 \mathcal L=A+B-C-D.
 }
\]
Substituting the explicit expressions for $A$, $B$, $C$, and $D$ gives
\[
 \boxed{
 \begin{aligned}
 \mathcal K={}&4\Id_2
 -s(W_1+W_2+W_3+W_4)\\
 &+s^2\bigl[W_4S+SW_1\bigr]
 -s^3W_4SW_1,
 \end{aligned}
 }
\]
y
\[
 \boxed{
 \begin{aligned}
 \mathcal L={}&s(W_1-W_2-W_3+W_4)\\
 &+s^2\bigl[W_4S-SW_1\bigr]
 +s^3W_4SW_1.
 \end{aligned}
 }
\]
The matrix $\mathcal K$ coincides with $R'(\mu)$ from (17), evaluated at the root. In
contrast, $\mathcal L$ measures how the residual changes when the initial physical state
is perturbed while the multiplier is momentarily held fixed. All matrices in these two
expressions are $2\times2$.

No derivatives of $W_i$ appear in these formulas. This is because
$\mathcal J_{\mathrm{ABBA}}$ is already the complete first derivative of the ABBA
composition: each $W_i=D_zf(z_i,t_i)$ is evaluated at its corresponding stage point.
Differentiating the $W_i$ again would produce second derivatives of the map and is not
part of the Jacobian sought here.

\subsection{Implicit derivative of the multiplier}

The converged multiplier satisfies
\[
 \mathcal R\bigl(z,\mu_h(z)\bigr)=0.
\]
Differentiating this identity with respect to $z$ and applying the chain rule,
\[
 D_z\mathcal R
 +D_\mu\mathcal R\,D\mu_h(z)=0.
\]
With the notation above,
\[
 \mathcal L+\mathcal K D\mu_h(z)=0.
\]
If $\mathcal K$ is invertible, the implicit function theorem locally guarantees that
$\mu_h$ is differentiable and
\[
 \boxed{
 D\mu_h(z)=-\mathcal K^{-1}\mathcal L.
 }
\]
This term is exactly the contribution that would be lost if $\mu$ were treated as a
constant independent of $z$ when differentiating the step.

\subsection{Differentiation of the final physical state}

We use the first of the two equivalent expressions in (13):
\[
 z_{n+1}=u^{(2)}+\mu_h(z_n).
\]
Como
\[
 u^{(0)}=z_n+\mu_h(z_n),
 \qquad
 v^{(0)}=z_n-\mu_h(z_n),
\]
their derivatives with respect to $z_n$ are
\[
 \frac{\partial u^{(0)}}{\partial z_n}
 =\Id_2+D\mu_h(z_n),
 \qquad
 \frac{\partial v^{(0)}}{\partial z_n}
 =\Id_2-D\mu_h(z_n).
\]
To make the chain rule explicit, we introduce the notation
\[
 M_\mu:=D\mu_h(z_n),
 \qquad
 U_j:=\frac{\partial u^{(j)}}{\partial z_n},
 \qquad
 V_j:=\frac{\partial v^{(j)}}{\partial z_n}.
\]
Thus, the two initial derivatives are
\[
 U_0=\Id_2+M_\mu,
 \qquad
 V_0=\Id_2-M_\mu.
\]

We now differentiate the four ABBA stages directly with respect to $z_n$. In the first stage,
\[
 u^{(1)}=u^{(0)}+s f(v^{(0)},t_n),
\]
and the chain rule gives
\[
 U_1=U_0+sW_1V_0.
\]
For the two central stages,
\[
 v^{(1)}=v^{(0)}+s f(u^{(1)},t_n),
 \qquad
 v^{(2)}=v^{(1)}+s f(u^{(1)},t_{n+1}),
\]
we obtain
\begin{align*}
 V_1&=V_0+sW_2U_1,\\
 V_2&=V_1+sW_3U_1
     =V_0+s(W_2+W_3)U_1
     =V_0+sS U_1.
\end{align*}
Finally, the last stage is
\[
 u^{(2)}=u^{(1)}+s f(v^{(2)},t_{n+1}),
\]
so that
\[
 U_2=U_1+sW_4V_2.
\]

We now substitute the expressions for $U_1$ and $V_2$ into that for $U_2$, without omitting
any steps:
\begin{align*}
 U_2
 &=U_1+sW_4\bigl(V_0+sS U_1\bigr)\\
 &=\bigl(\Id_2+s^2W_4S\bigr)U_1+sW_4V_0\\
 &=\bigl(\Id_2+s^2W_4S\bigr)
   \bigl(U_0+sW_1V_0\bigr)+sW_4V_0\\
 &=\bigl(\Id_2+s^2W_4S\bigr)U_0
   +\Bigl[sW_1+sW_4+s^3W_4SW_1\Bigr]V_0\\
 &=\underbrace{\bigl(\Id_2+s^2W_4S\bigr)}_{A}U_0
   +\underbrace{\Bigl[s(W_1+W_4)+s^3W_4SW_1\Bigr]}_{B}V_0.
\end{align*}
The blocks computed previously appear here exactly,
\[
 A=\frac{\partial u^{(2)}}{\partial u^{(0)}},
 \qquad
 B=\frac{\partial u^{(2)}}{\partial v^{(0)}}.
\]
Therefore, substituting $U_0=\Id_2+M_\mu$ and $V_0=\Id_2-M_\mu$ explicitly gives
\begin{align*}
 \frac{\partial u^{(2)}}{\partial z_n}
 &=A\bigl(\Id_2+M_\mu\bigr)
   +B\bigl(\Id_2-M_\mu\bigr)\\
 &=A+AM_\mu+B-BM_\mu\\
 &=A+B+(A-B)M_\mu\\
 &=A+B+(A-B)D\mu_h(z_n).
\end{align*}
If the expressions for $A$ and $B$ are also substituted, the unabbreviated result is
\begin{align*}
 \frac{\partial u^{(2)}}{\partial z_n}
 ={}&\Id_2+s(W_1+W_4)+s^2W_4S+s^3W_4SW_1\\
 &+\Bigl[
 \Id_2-s(W_1+W_4)+s^2W_4S-s^3W_4SW_1
 \Bigr]D\mu_h(z_n).
\end{align*}
The order of the products matters: $D\mu_h(z_n)$ appears on the right because Jacobians
act on perturbations as matrices multiplying from the left.

Differentiating the final correction $+\mu_h(z_n)$ as well gives
\[
 \boxed{
 D\Psi_{h,t_n}(z_n)
 =A+B+(A-B+\Id_2)D\mu_h(z_n).
 }
\]
Finally, substituting the implicit derivative of the multiplier,
\[
 \boxed{
 D\Psi_{h,t_n}(z_n)
 =A+B-(A-B+\Id_2)\mathcal K^{-1}\mathcal L.
 }
\]

To abbreviate this expression, we define
\[
 \mathcal P:=A+B,
 \qquad
 \mathcal Q:=A-B+\Id_2.
\]
Their explicit values are
\[
 \boxed{
 \mathcal P
 =\Id_2+s(W_1+W_4)+s^2W_4S+s^3W_4SW_1,
 }
\]
\[
 \boxed{
 \mathcal Q
 =2\Id_2-s(W_1+W_4)+s^2W_4S-s^3W_4SW_1.
 }
\]
Consequently, the exact Jacobian of the implicit physical step is given by the compact formula
\[
 \boxed{
 D\Psi_{h,t_n}(z_n)
 =\mathcal P-\mathcal Q\mathcal K^{-1}\mathcal L.
 }
\]
This formula includes both the direct variation of the four ABBA stages and the indirect
variation due to the dependence $\mu_h=\mu_h(z_n)$.

\subsection{Entrywise expression}

Since all the preceding matrices are $2\times2$, the only required inverse can be written
in closed form. Let
\[
 \mathcal P=(p_{ij}),
 \qquad
 \mathcal Q=(q_{ij}),
 \qquad
 \mathcal K=
 \begin{pmatrix}k_{11}&k_{12}\\k_{21}&k_{22}\end{pmatrix},
 \qquad
 \mathcal L=
 \begin{pmatrix}\ell_{11}&\ell_{12}\\\ell_{21}&\ell_{22}\end{pmatrix}.
\]
Si
\[
 \Delta_{\mathcal K}
 :=\det\mathcal K
 =k_{11}k_{22}-k_{12}k_{21}\neq0,
\]
then
\[
 \mathcal K^{-1}\mathcal L
 =\frac{1}{\Delta_{\mathcal K}}
 \begin{pmatrix}
  k_{22}\ell_{11}-k_{12}\ell_{21}
  &k_{22}\ell_{12}-k_{12}\ell_{22}\\
  -k_{21}\ell_{11}+k_{11}\ell_{21}
  &-k_{21}\ell_{12}+k_{11}\ell_{22}
 \end{pmatrix}.
\]
Let
\[
 \mathcal T:=\mathcal K^{-1}\mathcal L
 =\begin{pmatrix}\tau_{11}&\tau_{12}\\\tau_{21}&\tau_{22}\end{pmatrix}.
\]
The physical Jacobian is then
\[
 \boxed{
 D\Psi_{h,t_n}=
 \begin{pmatrix}
  p_{11}-q_{11}\tau_{11}-q_{12}\tau_{21}
  &p_{12}-q_{11}\tau_{12}-q_{12}\tau_{22}\\[1mm]
  p_{21}-q_{21}\tau_{11}-q_{22}\tau_{21}
  &p_{22}-q_{21}\tau_{12}-q_{22}\tau_{22}
 \end{pmatrix}.
 }
\]
That is,
\[
 D\Psi_{h,t_n}(z_n)=
 \begin{pmatrix}
  \dfrac{\partial X_{n+1}}{\partial X_n}
  &\dfrac{\partial X_{n+1}}{\partial Y_n}\\[2mm]
  \dfrac{\partial Y_{n+1}}{\partial X_n}
  &\dfrac{\partial Y_{n+1}}{\partial Y_n}
 \end{pmatrix}.
\]

As a check, for $h=0$ we have $s=0$ and therefore
\[
 \mathcal K=4\Id_2,
 \qquad
 \mathcal L=0,
 \qquad
 \mathcal P=\Id_2,
 \qquad
 \mathcal Q=2\Id_2.
\]
The preceding formula gives $D\Psi_{0,t_n}=\Id_2$, as required for the identity map.

It is this physical matrix, rather than $R'(\mu)=\mathcal K$, that must be used to verify directly
\[
 (D\Psi_{h,t_n})^TJ_0D\Psi_{h,t_n}=J_0,
\]
or, equivalently in two dimensions,
\[
 \det D\Psi_{h,t_n}=1.
\]

\section{Symmetry of the complete projected method}

The equation defining the projected step can be written as
\[
 \boxed{
 \Y_{n+1}-G^T\mu
 =\Phi^{\mathrm{ABBA}}_{h,t_n}
  \bigl(\Y_n+G^T\mu\bigr),
 \qquad
 \Y_n,\Y_{n+1}\in\M.
 } \tag{19}
\]
Applying the inverse of ABBA and using its symmetry,
\[
 \Y_n+G^T\mu
 =\Phi^{\mathrm{ABBA}}_{-h,t_{n+1}}
  \bigl(\Y_{n+1}-G^T\mu\bigr).
\]
We define the multiplier of the reverse step as
\[
 \mu_{\mathrm{inv}}=-\mu.
\]
Then
\[
 \Y_n-G^T\mu_{\mathrm{inv}}
 =\Phi^{\mathrm{ABBA}}_{-h,t_{n+1}}
  \bigl(\Y_{n+1}+G^T\mu_{\mathrm{inv}}\bigr),
\]
which is exactly equation (19) applied backward. Therefore, if the root is locally unique
and is solved exactly,
\[
 \boxed{
 \Psi_{-h,t_{n+1}}\circ\Psi_{h,t_n}=\Id.
 } \tag{20}
\]
Here $\Psi$ denotes ABBA together with Hairer's projection. In an implementation, the
iteration is stopped at a finite tolerance; reversibility is preserved up to that tolerance.

\section{Symplecticity of the complete projected method}

We now prove that the physical map obtained after symmetric projection is symplectic.
The point requiring care is that the multiplier is not constant: although it is found
iteratively within each time step, its converged value depends implicitly on the initial
physical state $z_n$.

\subsection{Diagonal and correction directions}

Recall the matrices introduced above
\[
 E=
 \begin{pmatrix}
  \Id_2\\ \Id_2
 \end{pmatrix},
 \qquad
 N=G^T=
 \begin{pmatrix}
  \Id_2\\ -\Id_2
 \end{pmatrix}.
\]
Thus, $Ez=(z,z)$ belongs to the diagonal manifold, while
$N\mu=(\mu,-\mu)$ gives the correction direction. If
\[
 z_{n+1}=\Psi_{h,t_n}(z_n),
\]
then equation (19) is equivalent to
\[
 \boxed{
 \PhiAB\bigl(Ez+N\mu_h(z)\bigr)
 =E\Psi_{h,t_n}(z)-N\mu_h(z).
 } \tag{21}
\]
The same multiplier is used on both sides of this equation within one time step, with
opposite signs. This does not mean that the multiplier remains unchanged in the next step.

\subsection{The residual and the implicit dependence of the multiplier}

In (11)--(12), $z_n$ was fixed and the residual was therefore written only as a function
of $\mu$. To study derivatives of the complete physical map, we restore this dependence
and write
\[
 \boxed{
 \mathcal R(z,\mu)
 :=G\PhiAB(Ez+N\mu)+2\mu.
 } \tag{22}
\]

\subsubsection*{Detailed construction of equation (22)}

Expression (22) combines in a single function the two conditions required for the two
state copies to coincide again after applying ABBA and the symmetric correction. We now
reconstruct it step by step.

We start from a physical state $z\in\R^2$. The matrix $E$ embeds it in the diagonal
manifold of the duplicated space:
\[
 Ez=
 \begin{pmatrix}\Id_2\\ \Id_2\end{pmatrix}z
 =\begin{pmatrix}z\\z\end{pmatrix}
 =(z,z)\in\R^4.
\]
The matrix $N=G^T$ introduces the multiplier in the direction transverse to that manifold:
\[
 N\mu=
 \begin{pmatrix}\Id_2\\-\Id_2\end{pmatrix}\mu
 =\begin{pmatrix}\mu\\-\mu\end{pmatrix}
 =(\mu,-\mu)\in\R^4.
\]
Therefore, the ABBA argument in (22) is the pre-corrected state
\[
 Ez+N\mu=(z+\mu,z-\mu).
\]
The two copies are separated before integration by equal corrections with opposite signs.
This separation is intentional: $\mu$ is chosen so that the complete symmetric correction
makes both copies coincide again at the end of the step.

Applying the ABBA map gives an auxiliary state in the duplicated space:
\[
 \PhiAB(Ez+N\mu)
 =\begin{pmatrix}u^{(2)}(z,\mu)\\v^{(2)}(z,\mu)\end{pmatrix}.
\]
Here the dependence on $z$ and $\mu$ has been written explicitly. In general, the four
ABBA stages give
$u^{(2)}(z,\mu)\neq v^{(2)}(z,\mu)$.

The matrix
\[
 G=(\Id_2\; -\Id_2):\R^4\longrightarrow\R^2
\]
measures exactly this mismatch, since for every $(u,v)\in\R^4$,
\[
 G\begin{pmatrix}u\\v\end{pmatrix}=u-v.
\]
In particular,
\[
 G\PhiAB(Ez+N\mu)
 =u^{(2)}(z,\mu)-v^{(2)}(z,\mu).
\]
This term is the diagonal-constraint defect before applying the output correction.

The symmetric projection uses the same multiplier at the output. The final duplicated state is
\[
 \Y_{n+1}
 =\PhiAB(Ez+N\mu)+N\mu.
\]
To require this state to belong to $\M$, we impose
\[
 G\Y_{n+1}=0.
\]
Substituting the preceding expression,
\begin{align*}
 G\Y_{n+1}
 &=G\PhiAB(Ez+N\mu)+GN\mu.
\end{align*}
Since $N=G^T$ and $GG^T=2\Id_2$, we obtain
\[
 GN\mu=GG^T\mu=2\mu.
\]
This calculation explains the term $2\mu$ in (22): it is neither an approximation nor an
arbitrary factor, but the exact contribution of the output correction to the constraint
defect. Consequently,
\[
 G\Y_{n+1}
 =G\PhiAB(Ez+N\mu)+2\mu
 =\mathcal R(z,\mu).
\]
Therefore, the following statements are equivalent:
\[
 \boxed{
 \mathcal R(z,\mu)=0
 \quad\Longleftrightarrow\quad
 G\Y_{n+1}=0
 \quad\Longleftrightarrow\quad
 \Y_{n+1}\in\M.
 }
\]

Writing the equality componentwise, with
$u^{(2)}=(u_X^{(2)},u_Y^{(2)})^T$,
$v^{(2)}=(v_X^{(2)},v_Y^{(2)})^T$, and
$\mu=(\mu_X,\mu_Y)^T$, the nonlinear system (22) is
\[
 \mathcal R(z,\mu)=
 \begin{pmatrix}
  u_X^{(2)}(z,\mu)-v_X^{(2)}(z,\mu)+2\mu_X\\
  u_Y^{(2)}(z,\mu)-v_Y^{(2)}(z,\mu)+2\mu_Y
 \end{pmatrix}
 =\begin{pmatrix}0\\0\end{pmatrix}.
\]
Thus, two equations are solved for the two components of $\mu$. Four equations are not
needed: membership in the diagonal only requires the difference between the copies to vanish.

It is also useful to check the dimensions of each object:
\[
 z,\mu\in\R^2,
 \qquad
 E,N\in\R^{4\times2},
 \qquad
 Ez+N\mu\in\R^4,
\]
\[
 \PhiAB:\R^4\to\R^4,
 \qquad
 G\in\R^{2\times4},
 \qquad
 \mathcal R(z,\mu)\in\R^2.
\]
Thus, (22) correctly defines a map
\[
 \mathcal R:\R^2\times\R^2\longrightarrow\R^2.
\]

When $z$ is fixed during a time step, one may abbreviate
$R(\mu)=\mathcal R(z,\mu)$, as in (11)--(12). However, differentiating the complete
physical map requires retaining the dependence on $z$, because changing the initial state
changes all four ABBA stages and hence the root of the projection system.

As a simple check, if $h=0$, then $\PhiAB=\Id_4$. In that case,
\[
 \mathcal R(z,\mu)
 =G(Ez+N\mu)+2\mu
 =GEz+GN\mu+2\mu.
\]
Because $GE=0$ and $GN=2\Id_2$, it follows that
\[
 \mathcal R(z,\mu)=4\mu.
\]
The unique root is $\mu=0$, as required when no time advance occurs. Moreover,
$D_\mu\mathcal R=4\Id_2$ for $h=0$; by continuity, this matrix generally remains
invertible for sufficiently small steps, allowing local use of the implicit function theorem.

The equation determining the multiplier is
\[
 \mathcal R(z,\mu)=0.
\]
Consequently, the multiplier obtained after convergence can be written locally as
\[
 \mu=\mu_h(z).
\]
This does not assert that an explicit formula for $\mu_h$ is available. It only states that
changing $z$ changes the nonlinear system and therefore its converged solution.

The Jacobian of the residual with respect to the multiplier is
\[
 D_\mu\mathcal R
 =G\,D\PhiAB\,N+2\Id_2. \tag{23}
\]
This is precisely the matrix $R'(\mu)$ obtained in (16). If it is invertible at the solution,
the implicit function theorem guarantees that $\mu_h(z)$ is locally differentiable.
In fact,
\[
 D\mu_h(z)
 =-\bigl(G\,D\PhiAB\,N+2\Id_2\bigr)^{-1}
   G\,D\PhiAB\,E,
\]
where $D\PhiAB$ is evaluated at $Ez+N\mu_h(z)$.

Define
\[
 M:=D\mu_h(z),
 \qquad
 J_\Psi:=D\Psi_{h,t_n}(z).
\]
Differentiating (21) with respect to $z$ gives
\[
 \boxed{
 D\PhiAB\,(E+NM)=EJ_\Psi-NM.
 } \tag{24}
\]
For a physical perturbation $u\in\R^2$, this identity states that
\[
 (E+NM)u
 \xrightarrow{\;D\PhiAB\;}
 (EJ_\Psi-NM)u.
\]
Here $Mu=D\mu_h(z)u$ compares converged multipliers associated with nearby initial states;
it is not the difference between two internal iterations.

\subsection{Symplectic orthogonality of the two directions}

Direct multiplication using the definition of $\Om$ gives
\[
 \boxed{
 \begin{aligned}
 E^T\Om E&=2J_0,
 &N^T\Om N&=-2J_0,\\
 E^T\Om N&=0,
 &N^T\Om E&=0.
 \end{aligned}
 } \tag{25}
\]
Therefore, the diagonal direction $(u,u)$ and correction direction $(a,-a)$ are
symplectically orthogonal with respect to $\Om$.

\subsection{Conservation of the extended symplectic form}

Equation (24) and the symplecticity of ABBA in (8) imply
\begin{align}
 &(EJ_\Psi-NM)^T\Om(EJ_\Psi-NM) \notag\\
 &\quad=\bigl[D\PhiAB(E+NM)\bigr]^T
          \Om
          \bigl[D\PhiAB(E+NM)\bigr] \notag\\
 &\quad=(E+NM)^T
          (D\PhiAB)^T\Om D\PhiAB
          (E+NM) \notag\\
 &\quad=(E+NM)^T\Om(E+NM). \tag{26}
\end{align}
The last equality is exactly where symplecticity of the duplicated ABBA map is used.
Thus, (26) states that $D\PhiAB$ preserves $\Om$ between the initial and final lifted perturbations.

Using the four identities in (25), the initial side becomes
\begin{align}
 (E+NM)^T\Om(E+NM)
 &=E^T\Om E+M^TN^T\Om NM \notag\\
 &=2J_0-2M^TJ_0M. \tag{27}
\end{align}
The mixed terms vanish because $E^T\Om N=N^T\Om E=0$. Likewise, the final side is
\begin{align}
 (EJ_\Psi-NM)^T\Om(EJ_\Psi-NM)
 &=J_\Psi^TE^T\Om EJ_\Psi+M^TN^T\Om NM \notag\\
 &=2J_\Psi^TJ_0J_\Psi-2M^TJ_0M. \tag{28}
\end{align}

Substituting (27) and (28) into (26) gives
\[
 2J_\Psi^TJ_0J_\Psi-2M^TJ_0M
 =2J_0-2M^TJ_0M. \tag{29}
\]
The term produced by the implicit dependence of the multiplier is identical on both sides
and cancels. Hence,
\[
 \boxed{J_\Psi^TJ_0J_\Psi=J_0.} \tag{30}
\]
Since $J_\Psi=D\Psi_{h,t_n}(z)$,
\[
 \boxed{
 D\Psi_{h,t_n}(z)^T J_0D\Psi_{h,t_n}(z)=J_0.
 } \tag{31}
\]
Therefore, the complete physical method, consisting of ABBA together with Hairer's
symmetric projection, is symplectic on the physical phase space. In two dimensions this also gives
\[
 \det D\Psi_{h,t_n}(z)=1,
\]
so the method preserves oriented area in the $(X,Y)$ plane.

The proof requires the following local hypotheses:
\begin{enumerate}
 \item the ABBA map is symplectic with respect to $\Om$;
 \item the equation $\mathcal R(z,\mu)=0$ has a locally unique solution;
 \item $D_\mu\mathcal R$ is invertible, so $\mu_h(z)$ is differentiable;
 \item the nonlinear projection equation is solved exactly.
\end{enumerate}
The reduced multiplier system defines the same exact projected map independently of the
particular convergent nonlinear solver. With a finite stopping tolerance, symplecticity is
preserved only up to the defect introduced by that tolerance.

\section{State extensions and dimensional conventions}

The \texttt{state\_extension} selector chooses which variables are duplicated.
For one guiding-center particle, the complete dimensional convention is:

\begin{center}
\footnotesize
\begin{tabular}{@{}>{\raggedright\arraybackslash}p{0.15\textwidth}
                    >{\raggedright\arraybackslash}p{0.18\textwidth}
                    >{\raggedright\arraybackslash}p{0.18\textwidth}
                    >{\raggedright\arraybackslash}p{0.16\textwidth}
                    >{\raggedright\arraybackslash}p{0.22\textwidth}@{}}
\toprule
Extension & Accepted state & Splitting state & Reduced unknown & Simultaneous unknown \\
\midrule
\texttt{physical}
  & $z\in\R^2$
  & $(u,v)\in\R^4$
  & $\mu\in\R^2$
  & $(u_f,v_f,\mu)\in\R^6$ \\
\texttt{shared\_time}
  & $(z,t,\kappa)\in\R^4$
  & $(u,v,t,k)\in\R^6$
  & $\mu\in\R^2$
  & $(u_f,v_f,\mu)\in\R^6$ \\
\texttt{fully\_extended}
  & $Z=(z,t,k)\in\R^4$
  & $(Z_1,Z_2)\in\R^8$
  & $\mu\in\R^4$
  & $(Z_{1f},Z_{2f},\mu)\in\R^{12}$ \\
\bottomrule
\end{tabular}
\end{center}

The literal $\R^2/\R^4/\R^6/\R^8/\R^{12}$ entries in the table are for one
particle.  For a physical run with $N$ independent particles, the accepted,
splitting, reduced, and simultaneous dimensions scale to $2N$, $4N$, $2N$,
and $6N$, respectively.  The shared-time and fully extended strategies
currently require $N=1$.

The accepted-state and observer-state dimensions are intentionally distinct
for \texttt{shared\_time}.  Physical and shared-time observers receive the
closed physical map $z\mapsto z_{n+1}$ in $\R^2$; the carried $t$ and
$\kappa$ histories remain in diagnostics.  Fully extended observers instead
receive the accepted internal map $Z\mapsto Z_{n+1}$ in $\R^4$.  The
diagnostics fields \texttt{observer\_state\_dimension} and
\texttt{observer\_state\_kind} state which convention applies.

The shared-time strategy duplicates only the physical state and shares one
$(t,k)$ pair.  Its accepted conjugate variable is normalized as $\kappa=k/2$.
Because the momentum update is triangular, it does not feed back into $z$ and
the physical trajectory agrees with the corresponding \texttt{physical}
configuration.  Its genuine $\R^6$ splitting state must not be confused with
the temporary $\R^6$ nonlinear unknown of the physical simultaneous
formulation.

The fully extended strategy duplicates the complete autonomous state, operates
on $(Z_1,Z_2)\in\R^8$, and advances direct $k$.  Its full-state projection can
define a different physical map.  In the simultaneous formulation, the two
four-component final copies and the four-component multiplier form the
$\R^{12}$ nonlinear workspace shown in the table.  Both non-physical state
extensions currently require exactly one particle so that these dimensions are
literal.  Every public ABBA class selects these strategies through
\texttt{state\_extension}; no separate extension method class is exposed.

\end{document}
